%% SpecTexNet manuscript (Elsevier elsarticle; works for EAAI, ESWA, Computers in Industry, etc.)
%% Numbers come from tables/results_macros.tex, written by scripts/aggregate.py.
%% Red "XX" means the experiment has not been run yet. Red \todo notes mark text that must be
%% checked against the real results before submission.
\documentclass[preprint,12pt,authoryear]{elsarticle}

\usepackage{amsmath,amssymb}
\usepackage{graphicx}
\usepackage{booktabs}
\usepackage{xcolor}
\usepackage{hyperref}
\usepackage{algorithm}
\usepackage{algpseudocode}
\usepackage{siunitx}

\newcommand{\todo}[1]{\textcolor{red}{[TODO: #1]}}
\newcommand{\figorbox}[2]{\IfFileExists{#1}{\includegraphics[width=#2]{#1}}{\fbox{\parbox{#2}{\centering\textcolor{red}{Run the pipeline to create \texttt{\detokenize{#1}}}}}}}
\newcommand{\tableor}[1]{\IfFileExists{#1}{\input{#1}}{\begin{center}\textcolor{red}{Table \texttt{\detokenize{#1}} not generated yet}\end{center}}}

\InputIfFileExists{tables/results_macros.tex}{}{}
\providecommand{\PropMacroFone}{\textcolor{red}{XX}}
\providecommand{\PropMacroFoneStd}{\textcolor{red}{XX}}
\providecommand{\PropMacroFoneCI}{\textcolor{red}{[XX, XX]}}
\providecommand{\PropBAcc}{\textcolor{red}{XX}}
\providecommand{\PropAcc}{\textcolor{red}{XX}}
\providecommand{\PropMCC}{\textcolor{red}{XX}}
\providecommand{\PropDDR}{\textcolor{red}{XX}}
\providecommand{\PropFAR}{\textcolor{red}{XX}}
\providecommand{\VoneMacroFone}{\textcolor{red}{XX}}
\providecommand{\VoneAcc}{\textcolor{red}{XX}}
\providecommand{\GainOverVone}{\textcolor{red}{XX}}
\providecommand{\NRaw}{\textcolor{red}{XX}}
\providecommand{\NClean}{\textcolor{red}{XX}}
\providecommand{\NDupGroups}{\textcolor{red}{XX}}
\providecommand{\NDupImages}{\textcolor{red}{XX}}
\providecommand{\NDerived}{\textcolor{red}{XX}}
\providecommand{\NNonPhoto}{\textcolor{red}{XX}}
\providecommand{\NNameConflicts}{\textcolor{red}{XX}}
\providecommand{\NCropped}{\textcolor{red}{XX}}
\providecommand{\ShortcutMetaFone}{\textcolor{red}{XX}}
\providecommand{\ShortcutPixFone}{\textcolor{red}{XX}}
\providecommand{\LeakyVoneAcc}{\textcolor{red}{XX}}
\providecommand{\LeakyVoneMacroFone}{\textcolor{red}{XX}}
\providecommand{\FriedmanP}{\textcolor{red}{XX}}
\providecommand{\NemenyiCD}{\textcolor{red}{XX}}

\journal{Engineering Applications of Artificial Intelligence}

\begin{document}
\begin{frontmatter}

\title{SpecTexNet: Spectral-residual guided attention and orientation-aware augmentation for fabric defect classification under a leakage-free evaluation protocol}

\author[bracu]{Author One\corref{cor1}}
\ead{author@example.com}
\author[bracu]{Author Two}
\cortext[cor1]{Corresponding author.}
\affiliation[bracu]{organization={Department of Computer Science and Engineering, BRAC University},
            city={Dhaka}, country={Bangladesh}}

\begin{abstract}
Automatic visual inspection of woven fabric has to recognise small, rare defects against a
strongly periodic background, and published results on public fabric datasets are hard to compare
because most use a single random split, report accuracy on heavily imbalanced data and do not check
for duplicate images. We address both the model and the evaluation. On the model side we propose
SpecTexNet, which adds a learnable multi-scale spectral-residual branch to a convolutional backbone.
The branch suppresses the regular weave pattern in the Fourier domain and passes the remaining
periodicity breaks to the backbone through a spectral-guided attention (SGA) module that gates
features both spatially and per channel. Because horizontal and vertical defects swap identity
under a quarter turn, we also introduce orientation-aware rotation augmentation (OA-Rot) and a
matching test-time augmentation that permute these labels instead of corrupting them. On the
evaluation side we audit the six-class Fabric Defects Dataset and find that \NDerived{} of its
\NRaw{} files are binarised copies of other files, that some files are masks rather than photographs,
and that defect-free and defective images come from different cameras, so that file metadata alone
reaches \ShortcutMetaFone\% macro-F1. We group every copy with its source, remove the remaining
duplicates with perceptual hashing, neutralise the camera shortcut, and run group-aware, stratified
5-fold cross-validation, reporting macro-F1, balanced accuracy and MCC with bootstrap
confidence intervals, McNemar tests and a Friedman/Nemenyi analysis. SpecTexNet with a ConvNeXt-T
backbone reaches a macro-F1 of \PropMacroFone{}\% $\pm$ \PropMacroFoneStd{} (95\% CI
\PropMacroFoneCI{}) and a balanced accuracy of \PropBAcc{}\%, compared with \VoneMacroFone{}\% for a
frozen-ResNet50 gradient-boosting pipeline representative of earlier work. Ablations isolate the
contribution of each component, and a lightweight EfficientNet-B0 variant runs on a single CPU
core in real time. Code, splits and trained models are public.
\end{abstract}

\begin{keyword}
Fabric defect detection \sep Textile inspection \sep Spectral residual \sep Attention \sep
Data augmentation \sep Data leakage \sep Statistical evaluation
\end{keyword}
\end{frontmatter}

%% ---------------------------------------------------------------------------
\section{Introduction}
\label{sec:intro}

Fabric inspection is still largely manual in many mills. A human inspector looking at fabric
moving at 15 to 20 m/min misses a substantial share of defects and tires quickly, and every
missed defect lowers the price a roll can be sold for \citep{kumar2008survey,ngan2011review}.
Camera-based inspection has been studied for three decades, first with hand-crafted texture
descriptors such as Fourier analysis \citep{chan2000fourier,tsai1999directional} and Gabor filter
banks \citep{kumar2002gabor}, and more recently with deep convolutional and transformer networks
\citep{kahraman2023review,carrilho2024survey,li2021survey}.

Two properties make fabric different from generic image classification. First, defect-free fabric
is close to periodic. Its Fourier spectrum is dominated by a few sharp peaks at the weave
frequencies, and a defect shows up as energy that does not fit that pattern. Classical methods
used this directly, while ImageNet-pretrained networks have to learn it from limited data. Second,
several defect classes are defined by orientation: a broken weft thread runs horizontally and a
broken warp thread runs vertically. This interacts badly with rotation augmentation, which is
otherwise one of the most useful augmentations for texture images.

There is also a measurement problem. Public fabric datasets are small, imbalanced and often
assembled from several sources. In the dataset used here, defect-free images make up roughly
60\% of the data, so a classifier that ignores defects altogether already reaches about 60\%
accuracy. Re-saved or resized copies of the same photograph can sit in both the training and the
test split, which inflates the score \citep{kaufman2012leakage}. In our own earlier pipeline
(frozen ResNet50 features with gradient-boosted trees), we found that ensemble weights had been fitted
on training data and that reported accuracy hid per-class F1 scores as low as 0.65 for the
minority classes. These issues are common, and they make it hard to know whether a new
architecture actually helps.

This paper makes four contributions.
\begin{enumerate}
\item \textbf{SpecTexNet}, a dual-branch classifier in which a learnable, multi-scale spectral
residual (LSR) branch computes where the weave periodicity is broken, and a spectral-guided
attention (SGA) module uses that map to gate backbone features. The branch adds about 0.4\,M
parameters and a few milliseconds of CPU time.
\item \textbf{Orientation-aware augmentation}. OA-Rot rotates training images by multiples of
$90^\circ$ and swaps the horizontal and vertical labels on odd turns. OA-TTA applies the same
permutation to predicted probabilities at test time.
\item \textbf{A dataset audit and leakage-free protocol} for the public six-class Fabric Defects
Dataset \citep{kaggle_fabric}. We show that binarised copies of the same photographs, masks and
camera-specific image formats let a model score well without detecting defects, group copies with
their sources, neutralise the camera shortcut, and evaluate with group-aware stratified folds,
imbalance-aware metrics and paired significance tests. We release the audit and the split file so
later work can compare on identical folds.
\item \textbf{A controlled comparison} of SpecTexNet with eight baselines, from frozen features with
gradient boosting to fine-tuned CNNs and vision transformers, all trained with the same recipe,
together with ablations, calibration, efficiency measurements, explanation maps and a zero-shot
cross-dataset test.
\end{enumerate}

%% ---------------------------------------------------------------------------
\section{Related work}
\label{sec:related}

\paragraph{Classical fabric inspection}
Surveys group classical methods into statistical, spectral, model-based and learning-based
approaches \citep{kumar2008survey,ngan2011review,hanbay2016review}. Spectral methods are the closest
to our work. \citet{chan2000fourier} used the Fourier spectrum of the fabric image to characterise
defects, and \citet{tsai1999directional} removed the dominant frequency components of directional
textures and reconstructed the image so that only anomalies remained. Gabor filter banks tuned to
the weave were used for the same purpose \citep{kumar2002gabor}. These methods are fast and
interpretable but rely on hand-set thresholds and tend to break when the weave, colour or lighting
changes.

\paragraph{Deep learning for fabric defects}
Recent reviews \citep{kahraman2023review,carrilho2024survey,li2021survey} describe a rapid shift to
CNN classifiers, object detectors and segmentation networks. Most classification studies fine-tune
an ImageNet backbone such as ResNet \citep{he2016resnet}, EfficientNet \citep{tan2019efficientnet}
or MobileNet \citep{howard2019mobilenetv3}, sometimes with added attention blocks
\citep{hu2018se,woo2018cbam}. Anomaly-detection benchmarks such as MVTec AD
\citep{bergmann2019mvtec} include textile-like categories but treat defects as a single anomaly
class. \todo{Add 5 to 8 fabric-specific papers from 2023 to 2026, preferably from the target
journal, and state in one sentence each what they do and how they evaluate (split, metric).}

\paragraph{Frequency-domain neural networks}
Several architectures use the Fourier domain inside the network. Fast Fourier Convolution
\citep{chi2020ffc} and Global Filter Networks \citep{rao2021gfnet} mix information globally with
learnable spectral filters, and FcaNet \citep{qin2021fcanet} uses DCT components as channel
descriptors. These modules learn generic spectral filters. SpecTexNet instead builds on the
spectral residual of \citet{hou2007spectral}, a saliency model whose prior (``remove the smooth
part of the log spectrum'') matches the structure of periodic textures, and makes its smoothing
kernel learnable.

\paragraph{Symmetry and augmentation}
Rotation-equivariant networks \citep{cohen2016gcnn} build symmetry into the architecture. When some
labels are themselves orientation dependent, the classifier should be equivariant only up to a
permutation of those labels. OA-Rot and OA-TTA implement this with data rather than architecture
changes, so they work with any backbone.

\paragraph{Evaluation practice}
Data leakage through duplicates is a known problem in applied machine learning
\citep{kaufman2012leakage}. For comparing classifiers, \citet{dietterich1998} and
\citet{demsar2006} recommend paired tests over folds or over test examples rather than comparing
mean accuracies. Calibration \citep{guo2017calibration} matters in inspection because a confidence
threshold decides whether a roll is stopped.

%% ---------------------------------------------------------------------------
\section{Data and evaluation protocol}
\label{sec:data}

\subsection{Dataset}
We use the Fabric Defects Dataset \citep{kaggle_fabric} (CC BY 4.0), which contains images of
fabric in six classes: \emph{defect-free}, \emph{hole}, \emph{horizontal}, \emph{lines},
\emph{stain} and \emph{vertical}. Before any training we audited every file, because the dataset
is assembled from several sources and this is not documented on the dataset page.
Table~\ref{tab:dataset} summarises the result.

\tableor{tables/table_dataset.tex}

\paragraph{Derived copies}
In the defect folders, most photographs \texttt{N.jpg} are accompanied by three files
\texttt{N\_processed (k).jpg} that are thresholded, binarised versions of the same photograph.
There are also files ending in ``\texttt{ - Copy}'' or ``\texttt{(1)}''. In total \NDerived{} of the
\NRaw{} files are derived from another file. A random split puts a photograph in the training set
and its binarised twin in the test set, which is the main way the published accuracy of earlier
pipelines on this dataset, including our own, was inflated (Section~\ref{sec:protocol-effect}). We
detect these files by name, put each one in the same group as its source, and use them for
training only. They never appear in validation or test folds, which contain only
independent photographs.

\paragraph{Non-photographs and naming conflicts}
\NNonPhoto{} files are almost binary images (at most four grey levels) without a source photograph,
for example line masks in the \emph{vertical} folder. They are removed. \NNameConflicts{} files carry a
class name in their filename that disagrees with their folder (for instance files named
\texttt{line\_\ldots} inside the \emph{hole} folder). We keep the folder label, list these files in
the released audit, and \todo{report how many of them were relabelled after manual inspection}.

\paragraph{Remaining duplicates}
Each remaining image is described by a 64-bit DCT perceptual hash (pHash) and a 64-bit difference
hash (dHash). Two images are linked when their MD5 checksums match, or when both Hamming distances
are small (pHash $\le 4$ and dHash $\le 6$). Requiring both hashes keeps distinct images of plain,
similar-looking fabric from being merged. Connected components define duplicate groups; groups
whose members carry different labels are removed as label noise and one image per group is kept.
\NDupImages{} images fall into \NDupGroups{} such groups, leaving \NClean{} independent images.
\todo{Inspect 30 random pairs from \texttt{data/duplicate\_pairs.csv} by eye and report how many are
true duplicates.}

\paragraph{Source shortcuts}
Most defect-free images are $4096\times256$ greyscale line-scan strips, while most defect images
are $640\times360$ colour photographs or 12-megapixel phone pictures. A random forest that sees
only width, height, aspect ratio, colour mode, file size and the number of grey levels reaches
\ShortcutMetaFone\% macro-F1 on our folds, without looking at a single pixel of fabric. We
therefore (i) centre-crop images whose aspect ratio exceeds 2 to a square (\NCropped{} images), so
that strips are not recognisable by their distortion after resizing, and (ii) convert every image to
greyscale. After these steps, a random forest on global intensity statistics of the network input
reaches \ShortcutPixFone\% macro-F1. \todo{Interpret: the closer this is to chance (16.7\% for six
balanced classes), the less a model can rely on the source. The ablation \emph{abl\_colour}
measures how much colour alone inflates the CNN.}

\subsection{Cross-validation}
We use stratified group 5-fold cross-validation (StratifiedGroupKFold, seed 42), so that class
proportions are kept and members of a duplicate group never appear in two folds. Within the four
training folds, 10\% of the groups are held out for early stopping and checkpoint selection. The
test fold is used only once, after training. Every model, including the re-evaluated earlier
pipeline, uses the same fold file, which we release.

\subsection{Metrics and statistics}
Macro-F1 is the primary metric because it weights every defect type equally. We also report
balanced accuracy (mean per-class recall), Matthews correlation coefficient, macro one-vs-rest
AUROC, expected calibration error (ECE, 15 bins), and plain accuracy for comparison with earlier
work. For deployment we report two binary figures obtained by merging all defect classes: the
defect detection rate (recall on defective images) and the false alarm rate (defect-free images
flagged as defective). Results are given as mean $\pm$ standard deviation over the five test
folds. 95\% confidence intervals are obtained from 2000 stratified bootstrap resamples of the pooled
out-of-fold predictions. Pairs of models are compared with an exact McNemar test on the pooled
predictions \citep{mcnemar1947}, and all models jointly with a Friedman test and the Nemenyi
critical difference over folds \citep{demsar2006}.

%% ---------------------------------------------------------------------------
\section{Method}
\label{sec:method}

Figure~\ref{fig:arch} gives an overview. An RGB image $x\in\mathbb{R}^{3\times H\times W}$ is
processed by a CNN backbone that produces a feature map $F\in\mathbb{R}^{C\times h\times w}$, and
in parallel by the spectral branch, which produces $G\in\mathbb{R}^{c\times h\times w}$. SGA
combines the two before global pooling and classification.

\begin{figure}[t]
\centering
\figorbox{figures/architecture.pdf}{0.95\linewidth}
\caption{SpecTexNet. The learnable spectral residual (LSR) turns the grey-level image into
periodicity-break maps at two scales. A light CNN encodes them into $G$, which gates the backbone
features $F$ spatially and per channel (SGA). Pooled $F'$ and $G$ are concatenated and
classified.}
\label{fig:arch}
\end{figure}

\subsection{Learnable multi-scale spectral residual}
Let $g$ be the grey-level image resized to $s\times s$ and mean-centred. To avoid the wrap-around
discontinuity of the discrete Fourier transform, $g$ is reflect-padded by $s/8$ on each side. With
$\mathcal{F}$ the 2-D DFT,
\begin{align}
A &= \log\left|\mathcal{F}(g)\right|,\qquad \Phi=\angle\,\mathcal{F}(g),\\
R &= A - k_\theta * A,\\
S &= \mathcal{G}_\sigma * \left|\mathcal{F}^{-1}\!\left(e^{R + i\Phi}\right)\right|^2 ,
\end{align}
where the convolution with $k_\theta$ is applied to the centred spectrum and $\mathcal{G}_\sigma$ is
a Gaussian blur ($\sigma=1.5$). \citet{hou2007spectral} fix $k$ to a $3\times 3$ mean filter. We
parameterise it as $k_\theta=\mathrm{softmax}(\theta)$ over the nine taps, initialised at
$\theta=0$, which gives exactly the mean filter at the start of training. The softmax keeps the
kernel positive and normalised, so $R$ stays a residual in the original sense, while gradients from
the classification loss can change which neighbouring frequencies count as ``background''.

For periodic fabric, $A$ consists of sharp weave peaks on a smooth envelope. The smooth envelope is
removed by the subtraction, and because the peaks are also present in $k_\theta * A$ they are
flattened. After inversion, regions that repeat the weave cancel out, and regions that break it
(holes, broken threads, scratches) stay bright. Figure~\ref{fig:spectral} shows examples. The
residual is computed at $s\in\{64,128\}$, because the weave pitch relative to the image size varies
between datasets. The resulting maps are cropped back, upsampled to $112\times112$, log-compressed
and standardised per image, so that the network sees where the saliency is concentrated rather
than its absolute level. Stains, which change colour over a large area without breaking the weave,
produce little spectral residual, so this class has to be recognised mostly by the backbone.

\begin{figure}[t]
\centering
\figorbox{figures/spectral_examples.pdf}{0.9\linewidth}
\caption{One image per class, its centred log-amplitude spectrum and the fixed spectral-residual
saliency at two scales (the initialisation of LSR).}
\label{fig:spectral}
\end{figure}

\subsection{Spectral encoder}
The two saliency maps are encoded by four $3\times3$ convolution, batch-norm and SiLU blocks with
stride 2 (widths 32, 64, 128, 128), and the output is average-pooled to the spatial size of $F$.
The encoder has about 0.4\,M parameters.

\subsection{Spectral-guided attention}
SGA lets the spectral evidence decide which positions and channels of the backbone features to
amplify:
\begin{align}
a &= \sigma\!\left(W_a\left[\,\phi(F)\,;\,G\,\right]\right)\in\mathbb{R}^{1\times h\times w},\\
m &= \sigma\!\left(\mathrm{MLP}(\mathrm{GAP}(G))\right)\in\mathbb{R}^{C},\\
F' &= F\odot(1+a)\odot(1+m),
\end{align}
where $\phi$ is a $1\times1$ convolution that reduces $F$ to 128 channels, $W_a$ is a $1\times1$
convolution and GAP is global average pooling. The gates use the residual form $(1+\cdot)$, so
attention can only amplify features and never removes the backbone's own evidence. The last
layers of both gates are initialised with zero weights and bias $-2$, so that
$a\approx m\approx 0.12$ and the network starts close to the plain backbone. The classifier
receives $[\mathrm{GAP}(F');\mathrm{GAP}(G)]$ after layer normalisation and dropout.

\subsection{Orientation-aware augmentation and test-time augmentation}
Let $\pi$ be the permutation that swaps \emph{horizontal} and \emph{vertical} and leaves the other
classes unchanged. OA-Rot samples horizontal and vertical flips with probability 0.5 each, which
preserve every label, and a rotation by $k\cdot90^\circ$ with $k$ uniform in $\{0,1,2,3\}$. When $k$ is
odd, the label $y$ is replaced by $\pi(y)$. The \emph{lines} class in this dataset contains defects
in several directions, so it is left unchanged. At test time OA-TTA averages the softmax outputs
over the identity, both flips and the two quarter turns, applying $\pi$ to the probability vector of
each rotated view:
\begin{equation}
\hat p(x)=\frac15\Big[p(x)+p(x^{\leftrightarrow})+p(x^{\updownarrow})
+\pi\big(p(\mathrm{rot}_{90}x)\big)+\pi\big(p(\mathrm{rot}_{270}x)\big)\Big].
\end{equation}
Ordinary rotation augmentation without the label swap (``naive rotation'') is included in the
ablation to measure how much the label noise costs.

\subsection{Loss}
We use the class-balanced focal loss \citep{cui2019cb,lin2017focal} with class weights
$w_c\propto(1-\beta)/(1-\beta^{n_c})$, $\beta=0.999$, where $n_c$ is the number of training images of
class $c$, and focusing parameter $\gamma=2$.

%% ---------------------------------------------------------------------------
\section{Experimental setup}
\label{sec:setup}

\paragraph{Baselines}
All baselines are trained on the same folds. (i) The earlier pipeline: frozen ImageNet ResNet50
features (2048-d) with second-order gradient-boosted trees implemented from scratch, three
regularisation settings and a convex ensemble whose weights are fitted on the validation split; a
single boosted model and a logistic-regression linear probe on the same features. (ii) Fine-tuned
MobileNetV3-L \citep{howard2019mobilenetv3}, EfficientNet-B0 \citep{tan2019efficientnet}, ResNet50
\citep{he2016resnet}, ConvNeXt-T \citep{liu2022convnext}, Swin-T \citep{liu2021swin} and DeiT-S
\citep{touvron2021deit}. All fine-tuned baselines use exactly the recipe of SpecTexNet (loss,
OA-Rot, OA-TTA, schedule), so differences between them come from the architecture.

\paragraph{Training details}
Images are resized to $224\times224$. Besides OA-Rot we use random resized crops with scale
$[0.7,1]$, so that small holes are rarely cropped out, mild colour jitter and Gaussian blur. We train
for up to 40 epochs with AdamW \citep{loshchilov2019adamw} (weight decay 0.05), batch size 32, a
head learning rate of $10^{-3}$ and $10^{-4}$ for the pretrained backbone, two warm-up epochs
followed by cosine decay, gradient clipping at 1.0, mixed precision and an exponential moving average
of the weights (decay 0.998). The checkpoint with the best validation macro-F1 is kept, with
patience 10. The implementation uses PyTorch \citep{paszke2019pytorch} and timm
\citep{wightman2019timm}. \todo{Name the GPU and report total training time from
\texttt{metrics.json}.}

%% ---------------------------------------------------------------------------
\section{Results}
\label{sec:results}

\subsection{Effect of the evaluation protocol}
\label{sec:protocol-effect}
To measure how much the evaluation protocol alone changes the numbers, we ran the earlier pipeline
(frozen ResNet50 features, scratch gradient boosting, weighted ensemble) twice: once on a split
that mimics the original study (random, ungrouped folds, binarised copies treated as independent
images, colour input, no cropping) and once under our protocol. Accuracy drops from
\LeakyVoneAcc\% to \VoneAcc\% and macro-F1 from \LeakyVoneMacroFone\% to \VoneMacroFone\%.
\todo{One or two sentences: this gap is the size of the leakage, and it is larger than most
architecture differences reported in the literature.}

\subsection{Comparison with baselines}
Table~\ref{tab:main} summarises the cross-validated results. SpecTexNet (ConvNeXt-T) reaches a
macro-F1 of \PropMacroFone{}\% and a balanced accuracy of \PropBAcc{}\%. The re-evaluated earlier
pipeline reaches \VoneMacroFone{}\% macro-F1 at \VoneAcc{}\% accuracy, a difference of
\GainOverVone{} percentage points in macro-F1. \todo{Describe the ranking in two or three sentences:
which family is strongest (fine-tuned CNN vs transformer), how large the gap between SpecTexNet and
its own backbone without the spectral branch is, and whether the accuracy column tells a different
story than macro-F1. Do not claim superiority where the McNemar p-value in
Table~\ref{tab:mcnemar} is above 0.05.}

\tableor{tables/table_main.tex}

The Friedman test over folds gives $p=\FriedmanP$ and the Nemenyi critical difference is
\NemenyiCD{} ranks (Figure~\ref{fig:cd}). \todo{State which models are not significantly different
from SpecTexNet. With five folds, rank tests have little power, so rely on the McNemar tests and
bootstrap intervals for pairwise claims.}

\begin{figure}[t]
\centering
\figorbox{figures/cd_diagram.pdf}{0.8\linewidth}
\caption{Average macro-F1 rank over folds (1 is best). Models closer than the critical difference
(CD) are not significantly different at $\alpha=0.05$.}
\label{fig:cd}
\end{figure}

\tableor{tables/table_mcnemar.tex}

\subsection{Per-class behaviour}
Table~\ref{tab:perclass} and Figure~\ref{fig:perclass} break the results down by class, and
Figure~\ref{fig:cm} shows the pooled confusion matrix of SpecTexNet. \todo{Discuss the classes that
were weakest in the earlier pipeline (hole, horizontal, vertical) and whether the spectral branch
helps them most, as Section 4.1 predicts. Also check whether stain is the class the spectral branch
helps least, which the method section also predicts.}

\tableor{tables/table_perclass.tex}

\begin{figure}[t]
\centering
\figorbox{figures/perclass_f1.pdf}{0.95\linewidth}
\caption{Per-class F1 on pooled out-of-fold predictions.}
\label{fig:perclass}
\end{figure}

\begin{figure}[t]
\centering
\figorbox{figures/confusion_proposed.pdf}{0.48\linewidth}\hfill
\figorbox{figures/reliability_proposed.pdf}{0.42\linewidth}
\caption{Left: pooled confusion matrix of SpecTexNet (ConvNeXt-T). Right: reliability diagram.}
\label{fig:cm}
\end{figure}

\subsection{Ablation study}
Table~\ref{tab:ablation} removes or replaces one component at a time on the ConvNeXt-T model.
\todo{Comment on each row in one sentence. The two augmentation rows are the key test of OA-Rot:
if naive rotation is worse than flips only, label noise from rotation is real, and if OA-Rot beats
both, the label swap recovers the benefit of rotation.}

\tableor{tables/table_ablation.tex}

\subsection{Industrial operating point and calibration}
Merging all defect types, SpecTexNet detects \PropDDR{}\% of defective images at a false alarm rate
of \PropFAR{}\% on defect-free images. \todo{Report ECE and say whether temperature scaling on the
validation split is needed before choosing a threshold.}

\subsection{Efficiency}
Table~\ref{tab:efficiency} lists parameters, multiply-accumulate operations and single-core CPU
latency. The spectral branch adds about 0.4\,M parameters. Its FFTs are counted analytically,
because the FLOP counter ignores them. \todo{Quote the latency overhead of the spectral branch
(SpecTexNet vs. its own backbone) and the throughput of the EfficientNet-B0 variant in images per
second.}

\tableor{tables/table_efficiency.tex}

\subsection{What the network looks at}
Figure~\ref{fig:explain} shows Grad-CAM maps \citep{selvaraju2017gradcam}, the spectral saliency
and the SGA spatial gate for correctly classified test images. \todo{Describe whether the gate and
Grad-CAM concentrate on the defect, and include at least one failure case.}

\begin{figure}[t]
\centering
\figorbox{figures/explain_grid.pdf}{0.85\linewidth}
\caption{Input, Grad-CAM for the true class, spectral-residual saliency (128 px) and SGA spatial
gate.}
\label{fig:explain}
\end{figure}

\subsection{Cross-dataset generalisation}
\todo{Run \texttt{scripts/cross\_dataset.py} on at least one external dataset (for example TILDA
or AITEX \citep{silvestre2019aitex}) with the five fold models as an ensemble and no fine-tuning.
Report binary AUROC and the aligned multi-class macro-F1 in a small table. If the drop is large,
say so plainly; reviewers trust that more than a claim of robustness.}

%% ---------------------------------------------------------------------------
\section{Discussion}
\label{sec:discussion}

\paragraph{Why a spectral prior helps}
A pretrained CNN learns textures from ImageNet, where periodic backgrounds are rare. The spectral
residual gives it, at almost no cost, a map of where the periodic structure is broken, which is
close to the definition of a weaving defect. Because the map is normalised per image, its
information lies in how concentrated the saliency is, which is also what separates a small hole
from clean fabric. The prior is weakest when the fabric is not periodic (printed or random
textures) or when the defect changes colour rather than structure, as with stains.

\paragraph{Measurement matters}
Re-evaluating our own earlier pipeline on leakage-free folds, with weights fitted on validation data
and with macro-F1 as the metric, changed the picture more than any architectural change.
\todo{Quantify: earlier reported accuracy 93.25\% on a single split vs. \VoneAcc{}\% accuracy and
\VoneMacroFone{}\% macro-F1 under the new protocol.}

\paragraph{Limitations}
The study uses one primary dataset with six image-level classes, and after removing copies some
defect classes have only a few dozen independent photographs, so per-class estimates for those
classes have wide confidence intervals. It does not localise defects
with boxes or masks, so it does not replace detection models where the defect position must be
reported. The duplicate thresholds were chosen by inspection and could merge genuinely different
images of very plain fabric. OA-Rot assumes that the only orientation-dependent classes are
horizontal and vertical. For datasets with diagonal defect classes, the permutation would need to be
extended.

%% ---------------------------------------------------------------------------
\section{Conclusion}
\label{sec:conclusion}
SpecTexNet adds a learnable spectral-residual branch and spectral-guided attention to a standard
CNN, and orientation-aware augmentation makes rotation usable for direction-dependent defect
classes. Under a de-duplicated, group-aware cross-validation protocol with imbalance-aware metrics,
the model reaches \PropMacroFone{}\% macro-F1. The released splits, code and safetensors weights let
others reproduce the numbers on identical folds. Future work will extend the spectral branch to
defect localisation and evaluate it on additional fabric types.

\section*{Data and code availability}
Code, the fold file, trained weights and a live demo are available at
\todo{GitHub URL} and \url{https://huggingface.co/spaces/codewithyeasin/fabric-defect-detection}.

\section*{CRediT authorship contribution statement}
\todo{Fill in.}

\section*{Declaration of competing interest}
The authors declare no competing interests.

\bibliographystyle{elsarticle-harv}
\bibliography{references}

\end{document}
